\documentclass[11pt,twocolumn]{article}

\usepackage[margin=1in]{geometry}
\usepackage{amsmath}
\usepackage{amssymb}
\usepackage{graphicx}
\usepackage{booktabs}
\usepackage{hyperref}
\usepackage{cite}
\usepackage{microtype}
\usepackage{xcolor}
\usepackage{enumitem}

\hypersetup{
    colorlinks=true,
    linkcolor=blue!60!black,
    citecolor=blue!60!black,
    urlcolor=blue!60!black,
}

% ── Metadata ──────────────────────────────────────────────────────────────────
\title{\textbf{Towards Channel-Efficient, Cross-Subject sEMG Gesture Recognition}\\
       \large Progress Report — putEMG Prime Project}

\author{Chris Dollo\\
        \small \texttt{dollochrisdavid@gmail.com}}

\date{May 2026}

% ── Document ──────────────────────────────────────────────────────────────────
\begin{document}
\maketitle

% ── Abstract ──────────────────────────────────────────────────────────────────
\begin{abstract}
Surface electromyography (sEMG) enables non-invasive hand gesture recognition for
prosthetic control and human-computer interfaces. However, most published systems
are trained and evaluated on data from the same subject, limiting practical
deployability. This report summarises progress on the \emph{putEMG Prime} project,
which aims to build a generalizable, channel-efficient sEMG gesture recognition
system using the publicly available putEMG dataset~\cite{kaczmarek2019putemg}.
We present within-subject and cross-subject evaluations of three deep learning
architectures and establish a strong baseline that exceeds the original
paper's benchmark under both protocols. We then outline a research roadmap
toward practical deployment: hand-crafted feature-based sequence models,
data-driven electrode reduction from 24 to 8 channels, and real-time inference.
\end{abstract}

% ── 1. Introduction ───────────────────────────────────────────────────────────
\section{Introduction}

Hand gesture recognition from surface electromyography is a well-studied problem
with applications in prosthetics, rehabilitation robotics, and gesture-based
human-computer interaction. Despite decades of research, a persistent gap remains
between laboratory accuracy and real-world deployability. Two related challenges
stand out:

\begin{enumerate}[noitemsep]
  \item \textbf{Cross-subject generalization.} Models trained on one participant's
        EMG signals typically perform poorly on a new user due to inter-subject
        variation in electrode placement, muscle anatomy, and signal amplitude.
  \item \textbf{Electrode count.} High-density electrode arrays (24+ channels)
        improve accuracy but are impractical for wearable or prosthetic devices.
        Reducing to 8 representative channels would greatly simplify hardware.
\end{enumerate}

This project addresses both challenges systematically, starting from a faithful
reproduction of the putEMG paper's within-subject benchmark before progressing to
cross-subject evaluation, feature-based modelling, and electrode reduction.

% ── 2. Dataset ────────────────────────────────────────────────────────────────
\section{Dataset}

The \textbf{putEMG} dataset~\cite{kaczmarek2019putemg} was collected at Poznan
University of Technology and is publicly available. Key properties are summarised
in Table~\ref{tab:dataset}.

\begin{table}[h]
\centering
\caption{putEMG dataset overview.}
\label{tab:dataset}
\begin{tabular}{ll}
\toprule
\textbf{Property} & \textbf{Value} \\
\midrule
Subjects          & 44 able-bodied participants \\
Electrodes        & 24 sEMG channels (forearm matrix) \\
Sampling rate     & 5120 Hz \\
Gesture classes   & 7 active (G1, G2, G3, G6, G7, G8, G9) \\
Repetitions       & $\sim$40 per gesture per subject \\
Sessions          & 2 per subject \\
Recording format  & CSV per session \\
\bottomrule
\end{tabular}
\end{table}

Each session contains multiple labelled gesture blocks. A MATLAB preprocessing
pipeline segments the blocks, extracts all 24 sEMG channels, and saves one
\texttt{combinedCell} matrix ($N_\text{reps} \times 7$) per subject. A subsequent
Python step bandpass-filters each channel (20--500 Hz, 4th-order Butterworth),
resamples each repetition to 1500 samples, and applies per-channel z-score
normalisation. The final per-subject file contains representations of shape
$(N, 1, 24, 1500)$ for deep learning, or $(N, 26, 3192)$ for feature-based models
(Section~\ref{sec:feature}).

% ── 3. Methods ────────────────────────────────────────────────────────────────
\section{Methods}

\subsection{Evaluation Protocols}

Two complementary evaluation protocols are used throughout the project.

\paragraph{Within-Subject 3-Fold CV.}
For each of the 44 subjects individually, repetitions are split into three
stratified folds. The model is trained on two folds and tested on the held-out
fold; this is rotated three times to yield three fold accuracies. The grand mean
and standard deviation are computed across all subjects. This protocol mirrors the
original putEMG paper's evaluation~\cite{kaczmarek2019putemg} and allows direct
comparison to their $\sim$90\% SVM/RMS benchmark.

\paragraph{Cross-Subject LOSO.}
Leave-One-Subject-Out (LOSO) cross-validation treats each of the 44 subjects as
a held-out test fold in turn. For a given test subject, the remaining 43 subjects'
data are pooled; 90\% is used for training and 10\% (stratified by class) is
reserved for early stopping only. No information from the test subject is seen
during training or validation. This protocol tests true generalization to unseen
individuals.

\subsection{Deep Learning Models}
\label{sec:dl}

All deep learning models receive raw sEMG windows of shape
$(B, 1, 24, 1500)$ and output class logits of shape $(B, 7)$. Three architectures
are compared.

\paragraph{EMG-TCN.}
A spatial mixing convolution ($C \times 1$ depthwise, projecting 24 channels to a
learned feature space) followed by four dilated temporal residual blocks. Each
residual block applies two dilated 1-D convolutions with dilation factors
$1, 2, 4, 8$ and a 1-D skip connection. Global average pooling and a linear
classifier head produce the final output. Dilation allows the network to capture
long-range temporal dependencies without excessive parameter counts.

\paragraph{EEGNet.}
A compact depthwise-separable CNN originally developed for EEG
classification~\cite{lawhern2018eegnet}. A temporal convolution learns
frequency-band filters, a depthwise convolution learns spatial filters per
channel, and a separable convolution fuses temporal and spatial features.
EEGNet serves as the standard BCI baseline due to its small parameter count
and strong published performance.

\paragraph{ShallowConvNet.}
Adapted from Schirrmeister et al.~\cite{schirrmeister2017deep}. A single temporal
convolution is followed by a spatial convolution across channels, a
square-then-log band-power nonlinearity, average pooling, and dropout. The
architecture approximates filter-bank common spatial patterns (FBCSP), making it
especially interpretable.

\paragraph{Training configuration.}
All models are trained with Adam ($\text{lr} = 10^{-3}$),
ReduceLROnPlateau scheduling (factor 0.5, patience 5), and early stopping
(patience 5, $\delta = 0.002$). Batch size is 16; maximum epochs 20; dropout
0.1. Best-epoch weights are saved per fold. Completed folds are skipped
automatically, making training resumable at any point without data loss.

\subsection{Feature-Based Models}
\label{sec:feature}

As an alternative to end-to-end learning, we extract 50 hand-crafted feature
groups from each channel using the \texttt{libemg} library. With 24 channels, each
window produces a 3192-dimensional feature vector
($50\ \text{groups} \times 24\ \text{channels}$). A sliding window of size 250
samples and stride 50 samples yields exactly 26 windows per 1500-sample
repetition. The resulting representation per repetition is a sequence matrix of
shape $(26, 3192)$.

Three sequence models are evaluated on this representation:

\begin{itemize}[noitemsep]
  \item \textbf{FeatureLSTM}: Bidirectional LSTM (2 layers, hidden 128); final
        hidden states from both directions concatenated and fed to a linear head.
        $\sim$3.8M parameters.
  \item \textbf{FeatureGRU}: Bidirectional GRU (2 layers, hidden 128); same
        classification head as FeatureLSTM. $\sim$2.8M parameters.
  \item \textbf{FeatureTransformer}: Sinusoidal positional encoding, 2-layer
        Transformer encoder ($d_\text{model}=128$, $n_\text{head}=4$,
        $d_\text{ff}=256$), mean pooling over the 26 timesteps, linear head.
        $\sim$675K parameters.
\end{itemize}

The feature cache is written once and shared between the within-subject and
cross-subject notebooks, avoiding redundant computation.

% ── 4. Results ────────────────────────────────────────────────────────────────
\section{Results}

\subsection{Within-Subject Evaluation}

Table~\ref{tab:within} reports grand mean accuracy across all evaluated subjects
(38 subjects; 6 subjects were excluded due to missing or incomplete data).

\begin{table}[h]
\centering
\caption{Within-subject 3-fold CV — deep learning (38 subjects).}
\label{tab:within}
\begin{tabular}{lcc}
\toprule
\textbf{Model} & \textbf{Mean Acc} & \textbf{Std} \\
\midrule
EMG-TCN        & \textbf{97.47\%}  & $\pm$2.78\% \\
ShallowConvNet & 93.50\%           & $\pm$4.09\% \\
EEGNet         & 91.05\%           & $\pm$5.56\% \\
\midrule
putEMG SVM+RMS~\cite{kaczmarek2019putemg} & $\sim$90\% & --- \\
\bottomrule
\end{tabular}
\end{table}

All three deep learning models exceed the published SVM+RMS benchmark. EMG-TCN
leads by more than 7 percentage points. EMG-TCN achieved 100\% on 5 individual
subjects (03, 12, 16, 26, 42), demonstrating that the task is near-perfectly
solvable with sufficient subject-specific data.

\subsection{Cross-Subject LOSO Evaluation}

Table~\ref{tab:loso} reports results after completing all 44 LOSO folds for all
three architectures.

\begin{table}[h]
\centering
\caption{Cross-subject LOSO — deep learning (44/44 folds).}
\label{tab:loso}
\begin{tabular}{lcccc}
\toprule
\textbf{Model} & \textbf{Mean} & \textbf{Std} & \textbf{Best} & \textbf{Worst} \\
\midrule
EMG-TCN        & \textbf{80.76\%} & 12.38\% & 99.29\% & 53.93\% \\
EEGNet         & 77.23\%          & 12.05\% & 98.21\% & 48.57\% \\
ShallowConvNet & 73.38\%          & 11.94\% & 95.00\% & 37.50\% \\
\bottomrule
\end{tabular}
\end{table}

EMG-TCN is the strongest cross-subject architecture, outperforming EEGNet by
3.53 pp and ShallowConvNet by 7.38 pp. The margin between models is consistent
and statistically meaningful given the full-population evaluation.

\paragraph{Key observations.}
\begin{enumerate}[noitemsep]
  \item \textbf{High variance.} All models exhibit $\sim$12 pp standard deviation.
        Some subjects are classified near-perfectly ($>$98\%) while others fall
        near chance ($<$55\%). This heterogeneity suggests that cross-subject
        generalization depends strongly on individual signal characteristics.
  \item \textbf{Val/test gap.} EMG-TCN achieves $\sim$93\% mean validation
        accuracy against $\sim$81\% test accuracy. Because the validation set is
        drawn from the same pool as training, this gap reflects inter-subject
        variance not captured during training rather than overfitting.
  \item \textbf{Consistent outliers.} Subject 06 is the hardest for both EMG-TCN
        (53.93\%) and EEGNet (48.57\%), while Subject 27 is the easiest across
        all three models (95--98\%). This consistency implies subject-level
        signal quality differences rather than random fold variation.
  \item \textbf{Cross-subject vs.\ within-subject gap.} The best cross-subject
        model (80.76\%) falls 16.7 pp below the best within-subject model
        (97.47\%), confirming that cross-subject generalization remains the
        primary bottleneck. Subject-adaptive fine-tuning is the most promising
        mitigation.
\end{enumerate}

% ── 5. Research Roadmap ───────────────────────────────────────────────────────
\section{Research Roadmap}

The project follows a seven-phase plan. Phases 1a and 2a are complete;
Phases 1b through 7 define the remaining objectives.

\paragraph{Phase 1b — Within-Subject Feature-Based Models.}
Evaluate FeatureLSTM, FeatureGRU, and FeatureTransformer under the same
3-fold within-subject protocol. The feature-based approach may offer better
interpretability and generalisation than end-to-end deep learning, especially
under limited per-subject data.

\paragraph{Phase 2b — Cross-Subject Feature-Based LOSO.}
Run full LOSO evaluation for all three sequence models on the 3192-dim feature
representation. Compare against the deep learning baselines from Phase 2a to
determine whether learned or hand-crafted features generalise better.

\paragraph{Phase 3 — Data-Driven Channel Reduction.}
Apply agglomerative hierarchical clustering to the 24 sEMG channels, grouping
correlated channels and selecting one representative per cluster. The target is
8 channels — one-third of the original count — sufficient for a realistic wrist
band or forearm sleeve. Clustering is performed separately for deep learning
(on raw signal correlations) and feature-based approaches (on feature-space
correlations).

\paragraph{Phase 4 — Anatomical Validation.}
Cross-reference the 8 selected channels with the electrode-to-muscle-group
mapping described in the putEMG paper. If the data-driven clusters align with
known functional muscle groups (e.g.\ flexor digitorum, extensor digitorum,
thenar), this provides an interpretable justification for the selection.

\paragraph{Phase 5 — Reduced-Channel Evaluation.}
Retrain all baseline models (both deep learning and feature-based) using only
the 8 representative channels. Compare accuracy against the 24-channel baselines
from Phases 1--2. The target is to approach within 5 pp of the full-channel
accuracy. Success here would demonstrate practical deployability.

\paragraph{Phase 6 — Per-Subject, Per-Class Analysis.}
Without any retraining, load each saved checkpoint and compute per-gesture
accuracy for every subject. Construct a subjects $\times$ gestures accuracy
matrix to identify which gesture classes are hardest to generalise across
subjects. This analysis motivates targeted data augmentation or gesture-specific
preprocessing in future work.

\paragraph{Phase 7 — Real-Time Data Collection (Future).}
Deploy the best-performing reduced-channel model in an online streaming setting.
Collect real-time sEMG data from new participants not present in the putEMG
dataset. Evaluate transfer performance and, if needed, study lightweight
per-user calibration.

% ── 6. Discussion ─────────────────────────────────────────────────────────────
\section{Discussion}

The within-subject results confirm that deep convolutional and temporal models
substantially outperform the SVM baseline from the original putEMG paper, even
without subject-specific hyperparameter tuning. The 97.47\% EMG-TCN accuracy
suggests that 7-class gesture recognition is effectively solved in the
within-subject setting with sufficient data.

The cross-subject results are more sobering. An 80.76\% mean LOSO accuracy with
a 12 pp standard deviation indicates that a single universal model struggles to
handle the full diversity of inter-subject signal variation. Several directions
may close this gap: domain adaptation, subject-specific batch normalisation
layers, or few-shot fine-tuning with a small number of calibration repetitions
from the new user.

The planned feature-based evaluation will clarify whether hand-crafted features
offer a complementary inductive bias. Features such as Mean Frequency and Slope
Sign Changes may capture aspects of muscle firing patterns that are more
invariant across subjects than raw waveform morphology.

% ── 7. Conclusion ─────────────────────────────────────────────────────────────
\section{Conclusion}

We have established strong within-subject and cross-subject deep learning
baselines for sEMG gesture recognition on the putEMG dataset. EMG-TCN achieves
97.47\% within-subject accuracy — exceeding the published benchmark — and 80.76\%
cross-subject LOSO accuracy. The gap between protocols quantifies the
cross-subject generalization challenge that motivates the remainder of the
project. The planned feature-based evaluation, channel reduction, and
anatomical validation phases will move the system progressively closer to
practical deployment in a real wearable device.

% ── References ────────────────────────────────────────────────────────────────
\bibliographystyle{ieeetr}
\begin{thebibliography}{9}

\bibitem{kaczmarek2019putemg}
P.~Kaczmarek, T.~Mankowski, and J.~Tomczynski,
``putEMG --- A Surface Electromyography Hand Gesture Recognition Dataset,''
\textit{Sensors}, vol.~19, no.~16, p.~3548, 2019.
\url{https://doi.org/10.3390/s19163548}

\bibitem{lawhern2018eegnet}
V.~J.~Lawhern, A.~J.~Solon, N.~R.~Waytowich, S.~M.~Gordon, C.~P.~Hung, and B.~J.~Lance,
``EEGNet: A Compact Convolutional Neural Network for EEG-Based Brain-Computer Interfaces,''
\textit{Journal of Neural Engineering}, vol.~15, no.~5, p.~056013, 2018.
\url{https://arxiv.org/abs/1611.08024}

\bibitem{schirrmeister2017deep}
R.~T.~Schirrmeister, J.~T.~Springenberg, L.~D.~J.~Fiederer, M.~Glasstetter,
K.~Eggensperger, M.~Tangermann, F.~Hutter, W.~Burgard, and T.~Ball,
``Deep Learning With Convolutional Neural Networks for EEG Decoding and Visualization,''
\textit{Human Brain Mapping}, vol.~38, no.~11, pp.~5391--5420, 2017.
\url{https://doi.org/10.1002/hbm.23730}

\end{thebibliography}

\end{document}
